\section{Task 2: Remaining Useful Life Regression}

\subsection{Exploration and preprocessing}
\label{sec:t2-prep}
% Numbers: kode/results/task2/eda.json (python -m task2_rul.eda). Figures: t2_*.pdf.

\paragraph{Data.} The training set holds 900 run-to-failure histories (133--778 cycles,
median 265); the test set holds 300 other machines whose logs stop at a random point
(50--455 cycles, median 163). Each cycle has three operating settings and 21 sensors.

\paragraph{Operating settings.} The settings take six discrete combinations
(Fig.~\ref{fig:t2-regimes}a); KMeans with $k=6$ on the standardised settings separates them
perfectly (largest distance to a centre: 0.03 standard deviations). The regime changes in
83\% of consecutive cycles and occurs equally often in train and test (16.5--17.0\% each).
It dominates the sensors: the regime explains 37--100\% of each sensor's variance
($\eta^2$), so raw readings jump between levels (Fig.~\ref{fig:t2-regimes}b) and hide the
slow wear trend.

\paragraph{Normalisation.} We z-score every sensor with the mean and standard deviation of
its regime, fitted on training machines only. Within-regime standard deviations are almost
identical across regimes (max/min ratio $\le 1.024$), so this removes the regime offset
without distorting the noise level. The degradation trend becomes visible
(Fig.~\ref{fig:t2-regimes}c), and the rank correlation with the capped RUL rises sharply,
e.g.\ for s04 from $|\rho|=0.28$ to $0.71$ (Fig.~\ref{fig:t2-sensors}a).

\begin{figure}[t]
  \centering
  \includegraphics[width=\linewidth]{t2_regimes}
  \caption{(a) Operating settings form six regimes. (b) Raw s04 of one machine, coloured by
  regime. (c) The same sensor after per-regime z-scoring shows the wear trend.}
  \label{fig:t2-regimes}
\end{figure}

\paragraph{Uninformative sensors.} For each training machine we compare the mean of its first
and last 20 cycles. Seven sensors (s01, s05, s06, s10, s16, s18, s19) show no change
(mean absolute change 0.23--0.25 $z$, the level expected from noise alone) and almost zero
correlation with RUL; the other 14 change by 1.22--2.66 $z$ (Fig.~\ref{fig:t2-sensors}b).
We drop the seven, using a threshold of 0.5 $z$ fitted on training machines only; fitting on
half of the machines selects the same sensors.

\begin{figure}[t]
  \centering
  \includegraphics[width=\linewidth]{t2_sensors}
  \caption{(a) $|$Spearman $\rho|$ between each sensor and the capped RUL, before and after
  per-regime normalisation. (b) Mean absolute change between a machine's first and last
  20 cycles; grey sensors are dropped.}
  \label{fig:t2-sensors}
\end{figure}

\paragraph{Missing values.} Each sensor misses about 2.0\% of its readings in both train and
test. Missingness is independent across sensors (6.61\% of rows miss two or more sensors,
6.69\% expected under independence), 98\% of gaps last a single cycle, and the rate does not
depend on the regime (2.02--2.04\%) or on the distance to failure (1.99\% in the last 25
cycles). The values therefore look missing completely at random. Because 35\% of rows miss
at least one sensor, dropping rows would discard too much; we forward-fill within each
machine (past values only) and set a reading with no earlier value to 0, the regime mean.

\paragraph{Spikes.} We compare each reading with the median of the previous 10 readings of
the same machine. The residual distribution (Fig.~\ref{fig:t2-spikes}b) has a Gaussian core
plus two side bumps at about $\pm 6$ noise standard deviations: isolated spikes of both signs
(51\% upward) with a median size of 3.8 $z$, on 0.2\% of readings in every sensor and in both
train and test. A reading more than 5 robust noise standard deviations from the trailing
median (the valley between core and bumps) is treated as missing and forward-filled
(Fig.~\ref{fig:t2-spikes}a). The filter only looks backwards, so the last cycle of a test
machine is cleaned exactly like a training cycle; we verified that preprocessing a history cut
at any cycle gives identical values to preprocessing the full history.

\begin{figure}[t]
  \centering
  \includegraphics[width=\linewidth]{t2_missing_spikes}
  \caption{(a) One machine before and after cleaning, with detected spikes and missing
  readings. (b) Residuals from the trailing median for all sensors (log scale); dashed lines
  mark the $\pm5$ noise-sd threshold.}
  \label{fig:t2-spikes}
\end{figure}

\paragraph{Two failure modes.} Clustering the per-machine change vectors gives two clear
groups (342 and 558 machines) in which different sensors degrade; s11 even drifts in
opposite directions (Fig.~\ref{fig:t2-modes}). Life lengths are similar (mean 287 vs 279
cycles). Early in life the modes look identical, so a model must infer the mode from the
trajectory itself; this favours models that combine several sensors nonlinearly over a single
linear trend per sensor.

\begin{figure}[t]
  \centering
  \includegraphics[width=\linewidth]{t2_failure_modes}
  \caption{Mean normalised reading against cycles before failure for the two failure modes.}
  \label{fig:t2-modes}
\end{figure}

\subsection{RUL labels}
% Numbers: kode/results/task2/labels.json (python -m task2_rul.labels).

Training machines run until they fail, so the last recorded cycle of each machine is its
failure cycle $T$, and the remaining life at cycle $t$ is $T-t$. Following the piecewise-linear
target of Heimes~\cite{heimes2008} and the exam's definition, we cap it:
\[
  \mathrm{RUL}(t) = \min\left(150,\; T - t\right).
\]
The cap reflects what the sensors can see. Early in life the normalised sensors are flat
(Fig.~\ref{fig:t2-modes}): a machine with 300 cycles left looks the same as one with 200 left, so
an uncapped label would ask the model to predict something that is not in the data. The wear
signal of an example machine starts to rise roughly where its capped label starts to fall
(Fig.~\ref{fig:t2-labels}a--b). The hidden test labels use the same cap.

Because 99\% of the machines live longer than 150 cycles, 47\% of all training rows carry the
label 150, while every value below the cap is shared by about 900 rows
(Fig.~\ref{fig:t2-labels}c). This imbalance matters later: evaluating on all rows would be
dominated by easy early-life cycles, which is one reason we validate on test-like cut points
(Section~\ref{sec:t2-validation}). Assertions in the code check that the label is 0 at every
failure cycle, lies in $[0,150]$, and (before capping) drops by exactly one per cycle.

\begin{figure}[t]
  \centering
  \includegraphics[width=\linewidth]{t2_labels}
  \caption{(a) Smoothed wear signal of one machine. (b) Its RUL label, uncapped and capped; the
  dotted line marks where the capped label starts to fall. (c) Distribution of all training
  labels.}
  \label{fig:t2-labels}
\end{figure}

\subsection{Sequence construction}
% Numbers: kode/results/task2/windows.json (python -m task2_rul.windows).

Following common practice for RUL regression~\cite{babu2016,zheng2017,li2018}, each sample is a
sliding window: the last $W$ cycles of one machine ending at cycle $t$, with target
$\mathrm{RUL}(t)$ (Fig.~\ref{fig:t2-windows}a). This gives every sample the same shape
($W \times 14$) whatever the length of the machine's history, and it reproduces the test
situation exactly: the input is the history up to a cycle, the output is the remaining life
at that cycle. For a test machine we use the single window ending at its last recorded cycle.

\paragraph{Different history lengths.} Histories range from 50 to 778 cycles
(Fig.~\ref{fig:t2-windows}b). A history shorter than $W$ is padded at the front by repeating its
first reading, the neutral assumption that nothing happened before cycle~1. Since the shortest
test history has 50 cycles, no test machine needs padding for $W \le 50$; we compare $W=30$ and
$W=50$ in Section~\ref{sec:t2-models}. With the default of starting training windows at $t=W$ this gives 228{,}074
($W=30$) or 210{,}074 ($W=50$) training windows. Rows are stored once per machine and windows are
cut out per batch, which keeps memory at about 16~MB instead of 365--561~MB.

\paragraph{No future leakage.} Three things keep future information out:
(i) all preprocessing statistics are fitted on training machines only and the preprocessing is
causal (Section~\ref{sec:t2-prep}); (ii) a window ending at $t$ contains only cycles $\le t$ of the same
machine; (iii) validation is split by machine (Section~\ref{sec:t2-validation}). Point (ii) is
checked in code on 2{,}000 random windows per setting: each window holds a single machine,
ends at its own cycle, has consecutive cycles, contains no later cycle, matches the
preprocessed rows, and carries $\mathrm{RUL}(t)$ as target. To confirm that this check can fail,
we shift every window one cycle into the future and verify that the check rejects it.

\begin{figure}[t]
  \centering
  \includegraphics[width=\linewidth]{t2_windows}
  \caption{(a) Three windows of $W=30$ cycles from one machine and their targets. (b) History
  lengths of training machines (full life) and test machines (cut before failure), with the
  window lengths we compare.}
  \label{fig:t2-windows}
\end{figure}

\subsection{Models}
\label{sec:t2-models}
% Numbers: kode/results/task2/baseline.json, evaluate.json.

\paragraph{How models are compared.} Before training any model we fix one evaluation protocol
(details in Section~\ref{sec:t2-validation}): 720 training machines fit the preprocessing and the
model, 180 held-out machines are cut at 5 random cycles each between 50 and their failure, and
the model predicts the RUL at each of these 900 cuts. The cut histories have almost the same
length distribution as the test histories (median 160 vs.\ 162 cycles; interquartile range
106--218 vs.\ 117--212).

\paragraph{Non-deep baseline.} A classical regressor cannot read a window directly, so we
summarise each window ($W=30$) by three standard features per sensor: the mean of the last 5
cycles (current level), the window mean (a steadier level) and the least-squares slope (rate of
change), 42 features in total. We fit a ridge regression (linear) and gradient-boosted trees
(300 iterations; the library's internal early stopping is disabled because it splits by row,
which would leak between rows of the same machine).

\begin{table}[t]
  \centering
  \caption{Baselines on the held-out machines at test-like cuts ($W=30$). Score: asymmetric
  score per prediction; "worst 5\%" is the share of the total score caused by the worst 5\% of
  predictions.}
  \label{tab:t2-baseline}
  \begin{tabular}{lrrrrr}
    \toprule
    Model & RMSE & Score mean & Score median & Worst 5\% & Train RMSE \\
    \midrule
    Ridge          & 23.3 & 37.2 & 1.94 & 78\% & 20.1 \\
    Boosted trees  & 21.9 & 62.0 & 1.21 & 90\% & 14.1 \\
    \bottomrule
  \end{tabular}
\end{table}

The trees are more accurate overall (Table~\ref{tab:t2-baseline}), especially close to failure
(Fig.~\ref{fig:t2-baseline}), yet their mean asymmetric score is worse. The score explains why:
a late error grows as $e^{d/10}$, so a prediction 70 cycles too late costs about 1{,}100
points while a typical prediction costs about 1. For the trees, the worst 5\% of predictions,
machines with true RUL 60--110 predicted near 150, make up 90\% of the total score. The mean
score therefore measures a few large late misses rather than typical accuracy, which we
address with uncertainty estimates in Section~\ref{sec:t2-uncertainty}. The gap between train
and held-out RMSE (14.1 vs.\ 21.9) shows that the trees partly memorise the training machines.

\begin{figure}[t]
  \centering
  \includegraphics[width=0.85\linewidth]{t2_baseline}
  \caption{Predicted vs.\ true RUL of the baselines at the test-like cuts of the held-out
  machines. Points above the diagonal are late (over-optimistic) predictions.}
  \label{fig:t2-baseline}
\end{figure}

\paragraph{Training protocol for the deep models.} All deep models predict $\mathrm{RUL}/150$
and are trained with mean squared error, Adam (learning rate $10^{-3}$, batch 512) and dropout.
From the 720 training machines, 72 are held aside for early stopping (10 test-like cuts each);
after every epoch we measure their RMSE, halve the learning rate after 3 epochs without
improvement, stop after 6, and restore the best epoch. The 180 held-out machines are never used
during training or model selection. Every model is trained with three seeds (0, 1, 2) and we
report mean $\pm$ standard deviation. Before training we overfit a single batch of 64 windows
with dropout switched off; every network reaches an RMSE below 0.1 cycles, which confirms that
the forward and backward passes work.

\paragraph{1D-CNN.} Following~\cite{babu2016,li2018}, three 1D convolution layers
(32, 64, 64 filters of width 5, ReLU) slide along the time axis of the window; each filter sees
all 14 sensors and learns a local pattern. The resulting feature map is flattened and passed
through a dense layer (64 units) to one output (156k parameters for $W=30$). The first version
overfits early: its training error keeps falling while the error on the early-stopping machines
is lowest after 1--5 epochs (Fig.~\ref{fig:t2-cnn}b). With 165k windows from only 648 machines,
consecutive windows overlap in 29 of 30 cycles, so one epoch shows every machine about 255
times. We therefore tested three standard remedies, one change at a time
(Table~\ref{tab:t2-cnn}).

\begin{table}[t]
  \centering
  \caption{1D-CNN variants, three seeds each. "Early stop" is the best RMSE on the 72 held-aside
  training machines (used for the choice); the other columns are on the 180 held-out machines.}
  \label{tab:t2-cnn}
  \begin{tabular}{lrrrr}
    \toprule
    Variant & Params & Early stop & RMSE & Score mean \\
    \midrule
    Original ($W=30$)                   & 156k & 22.35 & $22.93 \pm 0.60$ & $45 \pm 20$ \\
    Pooled head (global average)        &  37k & 24.50 & $24.11 \pm 0.82$ & $49 \pm 10$ \\
    Dropout 0.5 + weight decay          & 156k & 22.43 & $22.77 \pm 0.45$ & $30 \pm 5$ \\
    Longer window ($W=50$)              & 238k & 21.10 & $21.45 \pm 0.63$ & $32 \pm 8$ \\
    \bottomrule
  \end{tabular}
\end{table}

Averaging each filter over time (pooled head) cuts the parameters by 76\% but makes the model
worse: \emph{when} a pattern occurs in the window matters, since the most recent cycles say most
about the current state. Stronger regularisation leaves the RMSE unchanged but lowers and
stabilises the asymmetric score. A longer window helps most, because 50 cycles show the trend
more clearly than 30. We choose $W=50$ because it is best on the early-stopping machines; it
also beats both baselines on the held-out machines in RMSE (21.45 vs.\ 21.87 and 23.31) and in
mean score (32 vs.\ 62 and 37). It still reaches its best error after 1--2 epochs
(Fig.~\ref{fig:t2-cnn}b), and most remaining errors lie at true RUL 60--110, where wear has only
just started (Fig.~\ref{fig:t2-cnn}c).

\begin{figure}[t]
  \centering
  \includegraphics[width=\linewidth]{t2_cnn_W50}
  \caption{1D-CNN with $W=50$. (a) Overfitting a single batch. (b) Learning curves of three seeds
  (dashed: training, with dropout on; solid: early-stopping machines). (c) Predicted vs.\ true RUL
  on the held-out machines (seed 0).}
  \label{fig:t2-cnn}
\end{figure}

\paragraph{GRU.} As the second architecture we use a recurrent network, the other standard family
for RUL~\cite{zheng2017}. A gated recurrent unit (GRU)~\cite{cho2014} reads the window one cycle at
a time and keeps a hidden state of 64 numbers; at every cycle two gates decide how much of the old
state to keep and how much new information to write. We stack two GRU layers and pass the state
after the last cycle to the same dense head as the CNN (44.5k parameters). A GRU is a lighter
variant of the LSTM (two gates instead of three) with similar accuracy in practice. With the same
protocol and both window lengths, $W=50$ is again best on the early-stopping machines (20.15 vs.\
21.04). The GRU learns more slowly than the CNN (best epoch 3--6 instead of 1--2,
Fig.~\ref{fig:t2-gru}b) and overfits less.

\begin{table}[t]
  \centering
  \caption{All models on the 180 held-out machines at test-like cuts (deep models: mean $\pm$ std
  over three seeds; ensemble = average of the three seeds' predictions).}
  \label{tab:t2-models}
  \begin{tabular}{lrrrr}
    \toprule
    Model & Params & RMSE & Score mean & Ensemble RMSE \\
    \midrule
    Ridge (42 features)            & --   & 23.31            & 37.2         & -- \\
    Boosted trees (42 features)    & --   & 21.87            & 62.0         & -- \\
    1D-CNN, $W=50$                 & 238k & $21.45 \pm 0.63$ & $32 \pm 8$   & 20.89 \\
    GRU, $W=30$                    & 45k  & $21.59 \pm 0.10$ & $40 \pm 1$   & 21.45 \\
    GRU, $W=50$                    & 45k  & $\mathbf{20.70 \pm 0.10}$ & $34 \pm 6$ & \textbf{20.53} \\
    \bottomrule
  \end{tabular}
\end{table}

The GRU with $W=50$ is the most accurate model (Table~\ref{tab:t2-models}). Compared with the CNN it
needs five times fewer parameters and its results vary six times less between seeds. Reading the
window in order suits a slowly accumulating process such as wear. On the asymmetric score the
CNN and GRU are close (32 vs.\ 34), well within their seed-to-seed spread, and both beat the two
baselines. Averaging the three seeds lowers the error of every deep model slightly, which we use
for the uncertainty estimates in Section~\ref{sec:t2-uncertainty}.

\begin{figure}[t]
  \centering
  \includegraphics[width=\linewidth]{t2_gru_W50}
  \caption{GRU with $W=50$, same panels as Fig.~\ref{fig:t2-cnn}.}
  \label{fig:t2-gru}
\end{figure}

\paragraph{Why some predictions are far off.} % kode/results/task2/error_analysis.json
The largest errors of the CNN occur at true RUL 40--120, in both directions: for example a machine
predicted at 139 cycles with 45 left, and one predicted at 8 with 83 left. To understand them we
define a \emph{wear level} per cycle (the mean absolute deviation of the 14 normalised sensors
from the healthy start level, averaged over the previous 15 cycles, minus the machine's own level
in its first 40 cycles) and a machine's \emph{wear onset}, the number of cycles before failure at
which the wear level first exceeds 0.25. This uses the full history up to failure, so it serves
only for analysis, never as a model input. The onset varies widely between machines (median 80,
interquartile range 59--102 cycles). The two examples are the extremes
(Fig.~\ref{fig:t2-error}a): the first machine still looks new 45 cycles before failure and only
starts to wear 28 cycles before the end, while the second has been wearing slowly since 214
cycles before failure and already looks more worn than a typical machine at the moment it fails.
Across all held-out cuts with true RUL 40--120, the prediction error falls steadily with the wear
onset (Fig.~\ref{fig:t2-error}b): late wear onset leads to over-optimistic (late) predictions,
early onset to pessimistic ones. The correlation is between $-0.46$ and $-0.81$ for all three seeds
and onset thresholds of 0.15, 0.25 and 0.35, and the same machines cause the large errors for every
seed. The failure mode does not matter (mean absolute error 16.5 vs.\ 16.0 cycles for seed 0).
The sensors show how worn a machine is \emph{now}, not how fast it will continue to wear, so a
machine that has not started wearing at the prediction point cannot be told apart from a new one.
Part of the error is therefore irreducible for any window-based model, which motivates the
prediction intervals in Section~\ref{sec:t2-uncertainty}.

\begin{figure}[t]
  \centering
  \includegraphics[width=\linewidth]{t2_error_analysis}
  \caption{(a) Wear level of the CNN's worst late and worst early prediction and of a machine with
  median wear onset; dots mark the prediction points. (b) Prediction error against wear onset for
  held-out cuts with true RUL 40--120 (seed 0).}
  \label{fig:t2-error}
\end{figure}

\subsection{Validation}
\label{sec:t2-validation}
% Numbers: kode/results/task2/validation_by_range.json, validation_cv.json (python -m task2_rul.validation).

\paragraph{Protocol.} All splits are by machine, never by row: rows of one machine are nearly
identical neighbours, so a row split would let a model recognise machines instead of learning
wear. Preprocessing and sensor selection are refitted on the training machines of every split.
Because the hidden test asks for the RUL at a random stopping point, every held-out machine is
evaluated at 5 random cuts between cycle 50 and its failure; these cut histories match the test
lengths closely (Section~\ref{sec:t2-models}). Evaluating on all windows instead would be
dominated by easy early-life windows: for the baselines it lowers the RMSE from 21.9--23.3 to
19.3--20.7.

\paragraph{Error per RUL range.} Table~\ref{tab:t2-ranges} and Fig.~\ref{fig:t2-validation}a split
the 900 held-out cuts by true remaining life. Close to failure (0--25 cycles), where maintenance
decisions are made, the GRU is the most accurate model (RMSE 8.5). All models are worst at 51--100
cycles (RMSE about 30) and predict too much life there (mean error $+13$ to $+22$): these are the
machines whose wear starts late (Section~\ref{sec:t2-models}). Most of the asymmetric score comes
from the 26--50 range, where a late prediction is both likely and expensive. Healthy machines
(true RUL 150) are predicted too low by every model ($-10$ to $-16$), which is the safe direction.

\begin{table}[t]
  \centering
  \caption{RMSE (asymmetric score per prediction in brackets) per range of true RUL on the 180
  held-out machines; deep models: mean of three seeds, $W=50$.}
  \label{tab:t2-ranges}
  \small
  \begin{tabular}{lrrrrr}
    \toprule
    True RUL & 0--25 & 26--50 & 51--100 & 101--149 & 150 \\
    Cuts     & 122   & 116    & 227     & 172      & 263 \\
    \midrule
    Ridge         & 15.1 (4.4) & 23.4 (109) & 32.1 (75)  & 23.4 (13.3) & 16.3 (3.5) \\
    Boosted trees & 10.1 (1.8) & 19.2 (241) & 32.5 (110) & 22.7 (11.9) & 13.5 (2.6) \\
    1D-CNN        & 9.5 (1.6)  & 18.1 (115) & 29.7 (56)  & 20.2 (8.4)  & 18.7 (4.1) \\
    GRU           & 8.5 (1.2)  & 17.8 (126) & 29.6 (60)  & 20.7 (9.1)  & 15.9 (2.9) \\
    \bottomrule
  \end{tabular}
\end{table}

\paragraph{Five-fold cross-validation.} So far all numbers come from one set of 180 held-out
machines. To check that the ranking is not luck of that split, we divide the 900 machines into
five groups, hold out each group once, and refit everything (preprocessing, sensor selection,
models; deep models with seed 0) on the other four. The first fold is the split used above and
reproduces its numbers exactly; the same 14 sensors are selected in every fold. The GRU has the
lowest RMSE in all five folds (Fig.~\ref{fig:t2-validation}b): $19.20 \pm 0.83$ against
$20.51 \pm 0.93$ for the boosted trees, $21.12 \pm 1.02$ for the CNN and $21.99 \pm 0.87$ for
ridge regression, and also the lowest asymmetric score ($20 \pm 9$ per prediction vs.\ 27--34).
Across folds the CNN is not better than the boosted trees, so its advantage on the first split
was partly luck; the GRU's advantage is consistent.

\begin{figure}[t]
  \centering
  \includegraphics[width=\linewidth]{t2_validation}
  \caption{(a) RMSE per range of true RUL on the 180 held-out machines. (b) RMSE of each model in
  each fold of a 5-fold cross-validation by machine.}
  \label{fig:t2-validation}
\end{figure}

\subsection{Uncertainty and maintenance decisions}
\label{sec:t2-uncertainty}
% Numbers: kode/results/task2/uncertainty.json (python -m task2_rul.uncertainty).

\paragraph{Prediction intervals.} For the final model (GRU, $W=50$) we draw several predictions per
input in three standard ways: MC dropout~\cite{gal2016} (one network, dropout kept on, 50 passes),
a deep ensemble~\cite{lakshminarayanan2017} (five networks trained from different random starts,
seeds 0--4), and both combined (five networks $\times$ 20 passes). The mean of the samples is the
prediction and $\pm 1.645$ standard deviations the raw 90\% interval. On the 180 held-out machines
these raw intervals are far too narrow: they contain the true RUL in only 52\% (MC dropout), 20\%
(ensemble) and 63\% (combined) of the cuts (Table~\ref{tab:t2-unc}). The samples measure how unsure
the \emph{model} is, not how unpredictable the \emph{machine} is: a machine whose wear starts late
looks healthy to every network, so they agree on a wrong answer.

We therefore calibrate the intervals with split-conformal prediction: on the 72 machines held aside
from training (10 cuts each), we compute the ratio between the actual error and the predicted
standard deviation and take its 90\% quantile $q$; the calibrated interval is the prediction
$\pm\, q$ standard deviations. These machines also chose the best epoch, a mild reuse that we
accept because their weights were never fitted. After calibration all three methods cover 89--90\%
of the held-out cuts. The combined method gives the narrowest intervals (mean width 50 cycles),
because its standard deviation tracks the actual error best (rank correlation 0.45); we use it from
here on. Its intervals adapt to the situation (Fig.~\ref{fig:t2-unc}): close to failure they are
narrow (31 cycles wide, 95\% coverage), in mid-life they are wide (58 cycles). A constant-width
interval also reaches 90\% overall but is 49 cycles wide even close to failure. In the 51--100
range coverage stays at 74\% for every method: these are again the machines whose wear starts
late, for which even the uncertainty estimate is overconfident.

\begin{table}[t]
  \centering
  \caption{90\% prediction intervals of the GRU ($W=50$) on the 180 held-out machines: coverage
  (share of true RUL inside) and mean width in cycles.}
  \label{tab:t2-unc}
  \begin{tabular}{lrrrr}
    \toprule
    & \multicolumn{2}{c}{Raw} & \multicolumn{2}{c}{Calibrated} \\
    Method & Coverage & Width & Coverage & Width \\
    \midrule
    MC dropout (1 net, 50 passes)        & 52\% & 22 & 89\% & 54 \\
    Ensemble (5 nets)                    & 20\% & 10 & 89\% & 61 \\
    Ensemble + MC dropout (5 $\times$ 20) & 63\% & 27 & 90\% & 50 \\
    Constant width                       & --   & -- & 89\% & 54 \\
    \bottomrule
  \end{tabular}
\end{table}

\begin{figure}[t]
  \centering
  \includegraphics[width=\linewidth]{t2_uncertainty}
  \caption{(a) Coverage of the 90\% interval per range of true RUL (ensemble + MC dropout).
  (b) Prediction, calibrated 90\% interval and true RUL over the life of three held-out machines.}
  \label{fig:t2-unc}
\end{figure}

\paragraph{When to service a machine.} % kode/results/task2/maintenance.json
A maintenance decision is taken every cycle, so we replay each of the 180 held-out machines cycle
by cycle from cycle 50 until its failure and apply a rule at every cycle. We compare a
\emph{point rule} (service as soon as the predicted RUL is at most $L$) with a \emph{cautious rule}
(service as soon as the lower end of the calibrated 90\% interval is at most $L$). Assuming that
maintenance needs 10 cycles to arrange, a machine is serviced \emph{too late} if fewer than 10
cycles are left when the rule fires; the true number of cycles left at that moment is useful life
thrown away. Fig.~\ref{fig:t2-maint}a shows the trade-off for $L = 0, 2, \dots, 80$. To catch
every machine in time, the point rule needs $L=34$ and services machines on average 36.4 cycles
before failure; the cautious rule needs $L=14$ and wastes 38.1 cycles. The two curves overlap: in
our data the interval does not improve this decision. Near the decision point the interval is
about 40 cycles wide for almost every machine (standard deviation 5.7), and its width hardly
tracks the actual error (correlation 0.12), so ``lower end $\le 14$'' is nearly the same rule as
``prediction $\le 34$''. Our recommendation is therefore simple: service a machine when its
predicted RUL drops to about 35 cycles. With a 10-cycle lead time this caught all 180 held-out
machines in time, at the cost of about 36 cycles of unused life per machine
(Fig.~\ref{fig:t2-maint}b). The interval remains useful to communicate how reliable a single
prediction is, especially in mid-life, where it is wide. A shorter lead time or a different cost
of failure would move the threshold; the curves let the maintenance team choose.

\begin{figure}[t]
  \centering
  \includegraphics[width=\linewidth]{t2_maintenance}
  \caption{Replaying the 180 held-out machines. (a) Share of machines serviced too late (fewer than
  10 cycles left) against mean unused life, for thresholds $L=0,2,\dots,80$; large dots: cheapest
  setting with no machine too late. (b) True cycles left when service is triggered at those
  settings; dashed line: 10-cycle lead time.}
  \label{fig:t2-maint}
\end{figure}

\subsection{Final model and test predictions}
% kode/results/task2/final_shift.json, final_shift_simulation.json, final_train.json, test_predictions.json

The final model averages five GRUs ($W=50$, seeds 0--4), the best architecture in every
comparison. After validation we retrain it with the same recipe on all 900 training machines
(preprocessing and sensor selection on all 900; 810 machines for training, 90 for early stopping).

\paragraph{Score-aware shift.} Because the asymmetric score punishes late predictions more than
early ones, lowering every prediction by a few cycles lowers the score at some cost in RMSE.
Optimising the score alone would subtract 13 cycles, but on the held-out machines this raises the
RMSE from 20.5 to 24.1. Since both metrics are graded equally, we choose the shift by simulating
the test: 2{,}000 test sets of 300 validation machines (252 machines never used to fit the weights,
one random cut each, as in the real test). For each shift we divide the change of each metric by
its natural spread between test sets (RMSE: 1.14; summed score: 4{,}658), which puts both metrics
in the same unit, and take the shift with the largest net benefit. This is a shift of 3 cycles:
it lowers the summed score by 22\% (in all 2{,}000 simulated test sets) while the RMSE rises by
only 0.10 (still lower than without shift in 30\% of the sets). Shifts of 2--4 cycles are almost as
good; from 6 cycles on, the RMSE cost outweighs the score gain. The real test set is not used for
any choice.

\paragraph{Test predictions.} The submitted RUL is the mean prediction of the five GRUs at the last
recorded cycle of each test machine, minus 3 cycles, clipped to $[0,150]$ (median 104.8, range
3.1--146.5). Based on the five-fold cross-validation and the simulation we expect a test RMSE of
about 19--21 cycles; the asymmetric score varies strongly between test sets because a few late
predictions dominate it. The saved weights reproduce \texttt{prediksi\_rul.csv} exactly in about
15 seconds without retraining (\texttt{python -m task2\_rul.predict}).
